%% kapa: exam=2 round=31 year=2020 period=1 subject=appraisal_practice subject_ko=감정평가실무 seq=3 number=3 points=20
%% source: Q-Net 감정평가사 기출문제 (cst00309) artlSeq=5207673 file=raw/1008/5207673/1교시_감정평가실무.pdf
%% conversion: kordoc 4.24.51 -> tools/gen2.py (auto); review=proofread-gpt-astra; flags=none
\begin{longquestion}{【문제 3】}
감정평가사甲은 A지방법원 판사乙로부터 도시철도사업과 관련한 토지의 감정평가를 의뢰받았다. 감정평가사甲은 본 소송 과정에서 원고와 피고의 이해관계가 첨예하게 대립하고 있는 점을 확인하고 각자의 입장에서 대상 토지를 사전분석해 보기로 하였다. 관련 법규 및 이론을 참작하고 제시된 자료를 활용하여 다음의 물음에 답하시오. (20점) \par
\medskip\noindent 물음 1) 피고(사업시행자이자 매수인) 입장에서 주장할 것으로 판단되는 대상토지의 이용상황을 관련 법규 등을 근거로 검토한 후, 해당 이용상황에 따른 대상 토지를 감정평가 하시오. (10점) \par
\medskip\noindent 물음 2) 원고(피수용자이자 매도인) 입장에서 주장할 것으로 판단되는 대상토지의 이용상황을 관련 법규 등을 근거로 검토한 후, 해당 이용상황에 따른 대상 토지를 감정평가 하시오. (10점) \par
< 자료 1 > 사건 개요 \par
1. 평가의뢰인: A지방법원 판사乙 \par
2. 사건번호: 2020구합○○○○ 손실보상금 \par
3. 원고: 丙 \par
4. 피고: A시 \par
< 자료 2 > 기본적 사항 \par
1. 감정평가목적: 소송(감정목적물에 대한 수용 당시의 적정한 보상금 산정) \par
2. 감정목적물: A시 B구 C동 10-3번지 \par
3. 감정할 사항: 감정목적물에 대한 2019. 05. 19을 가격시점으로 한 적정한 시가(보상액) \par
4. 토지 변동내역 \par
\medskip\noindent (1) B동 10번지 \par
1) 2014. 05. 24: 건축허가 득함 \par
2) 2014. 12. 05: 분할되어 본번에 -1을 부함 \par
3) 2014. 12. 12: 건축물 사용승인 득함 \par
4) 2018. 05. 19: 분할되어 본번에 -2를 부함 \par
\medskip\noindent (2) B동 10-1번지 \par
1) 2014. 12. 05: B동 10번지에서 분할 \par
2) 2014. 12. 05: 지목변경 \par
3) 2018. 05. 19: 분할되어 본번에 -3을 부함 \par
5. 대로1류(폭20M～25M) 변동내역 \par
\medskip\noindent (1) 2008. 02. 15: A시 도시계획시설(도로) 결정 및 지형도면고시 \par
\medskip\noindent (2) 2016. 06. 09: A시 도시계획시설(도로) 결정(변경) 및 지형도면고시 \par
\medskip\noindent (3) 2016. 07. 15: 보상계획공고 \par
\medskip\noindent (4) 2016. 09. 30: 도시계획시설(도로) 실시계획인가고시 \par
\medskip\noindent (5) 2017. 10. 31: 사업준공완료 \par
6. 도시철도사업관련 \par
\medskip\noindent (1) 2018. 05. 24: 도시철도 A선 사업계획승인(「도시철도법」 제7조 제1항) \par
7. 감정평가 관련자료 \par
\medskip\noindent (1) 대상토지의 개요 \par
{\small\centering\renewcommand{\arraystretch}{1.25}
\resizebox{\ifdim\width>\linewidth\linewidth\else\width\fi}{!}{
\begin{tabular}{|c|c|c|c|c|}
\hline
\makecell{소재지} & \makecell{편입면적(m\textsuperscript{2})} & \makecell{지목} & \makecell{이용상황} & \makecell{공법상 제한사항} \\
\hline
\makecell{A시 B구 C동 10-3번지} & \makecell{19} & \makecell{도} & \makecell{도로} & \makecell{준주거지역,\\도시철도} \\
\hline
\end{tabular}}\par}

\medskip\noindent (2) 표준지 공시지가 \par
{\small\centering\renewcommand{\arraystretch}{1.25}
\resizebox{\ifdim\width>\linewidth\linewidth\else\width\fi}{!}{
\begin{tabular}{|c|c|c|c|c|c|c|c|c|c|}
\hline
\makecell{기호} & \makecell{소재지} & \makecell{면적(m\textsuperscript{2})} & \makecell{지목} & \makecell{이용\\상황} & \makecell{용도\\지역} & \makecell{도로\\교통} & \makecell{형상\\/지세} & \makecell{2016.01.01\\(원/m\textsuperscript{2})} & \makecell{2018.01.01\\(원/m\textsuperscript{2})} \\
\hline
\makecell{A} & \makecell{B구 C동\\7번지} & \makecell{500} & \makecell{대} & \makecell{단독\\주택} & \makecell{준주거} & \makecell{세로\\(가)} & \makecell{사다리\\/평지} & \makecell{630,000} & \makecell{750,000} \\
\hline
\makecell{B} & \makecell{B구 C동\\10번지} & \makecell{550} & \makecell{대} & \makecell{상업용} & \makecell{준주거} & \makecell{광대한면} & \makecell{세장형\\/평지} & \makecell{1,250,000} & \makecell{1,500,000} \\
\hline
\end{tabular}}\par}

※ 2018. 01. 01자 표준지 기호B는 도시계획시설(도시철도)에 30\% 저촉됨 \par
\medskip\noindent (3) 시점수정치 \par
1) 2018. 01. 01 ～ 2019. 05. 19: 1.09268 \par
2) 2016. 01. 01 ～ 2019. 05. 19: 1.15069 \par
\medskip\noindent (4) 개별요인 비교치 \par
1) 도로접면 \par
{\small\centering\renewcommand{\arraystretch}{1.25}
\resizebox{\ifdim\width>\linewidth\linewidth\else\width\fi}{!}{
\begin{tabular}{|c|c|c|c|c|c|}
\hline
\makecell{구분} & \makecell{광대한면} & \makecell{중로한면} & \makecell{소로한면} & \makecell{세로(가)} & \makecell{세로(불)} \\
\hline
\makecell{광대한면} & \makecell{1.00} & \makecell{0.91} & \makecell{0.85} & \makecell{0.80} & \makecell{0.72} \\
\hline
\makecell{중로한면} & \makecell{1.10} & \makecell{1.00} & \makecell{0.93} & \makecell{0.88} & \makecell{0.80} \\
\hline
\makecell{소로한면} & \makecell{1.18} & \makecell{1.07} & \makecell{1.00} & \makecell{0.94} & \makecell{0.86} \\
\hline
\makecell{세로(가)} & \makecell{1.25} & \makecell{1.14} & \makecell{1.06} & \makecell{1.00} & \makecell{0.91} \\
\hline
\makecell{세로(불)} & \makecell{1.38} & \makecell{1.25} & \makecell{1.17} & \makecell{1.10} & \makecell{1.00} \\
\hline
\end{tabular}}\par}

2) 형상 \par
{\small\centering\renewcommand{\arraystretch}{1.25}
\resizebox{\ifdim\width>\linewidth\linewidth\else\width\fi}{!}{
\begin{tabular}{|c|c|c|c|c|}
\hline
\makecell{구분} & \makecell{정방형} & \makecell{가장형} & \makecell{세장형} & \makecell{사다리} \\
\hline
\makecell{정방형} & \makecell{1.00} & \makecell{1.02} & \makecell{1.00} & \makecell{0.99} \\
\hline
\makecell{가장형} & \makecell{0.98} & \makecell{1.00} & \makecell{0.98} & \makecell{0.97} \\
\hline
\makecell{세장형} & \makecell{1.00} & \makecell{1.02} & \makecell{1.00} & \makecell{0.99} \\
\hline
\makecell{사다리} & \makecell{1.01} & \makecell{1.03} & \makecell{1.01} & \makecell{1.00} \\
\hline
\end{tabular}}\par}

3) 도시계획시설 \par
{\small\centering\renewcommand{\arraystretch}{1.25}
\resizebox{\ifdim\width>\linewidth\linewidth\else\width\fi}{!}{
\begin{tabular}{|c|c|c|c|}
\hline
\makecell{도시계획시설} & \makecell{일반} & \makecell{도로} & \makecell{도시철도} \\
\hline
\makecell{일반} & \makecell{1.00} & \makecell{0.93} & \makecell{0.85} \\
\hline
\makecell{도로} & \makecell{1.08} & \makecell{1.00} & \makecell{0.92} \\
\hline
\makecell{도시철도} & \makecell{1.18} & \makecell{1.09} & \makecell{1.00} \\
\hline
\end{tabular}}\par}

\medskip\noindent (5) 그 밖의 요인 보정치 \par
1) 표준지 기호A \par
① 2018. 01. 01 표준지 공시지가: 1.35 \par
② 2016. 01. 01 표준지 공시지가: 1.45 \par
2) 표준지 기호B \par
① 2018. 01. 01 표준지 공시지가: 1.50 \par
② 2016. 01. 01 표준지 공시지가: 1.65 \par
\medskip\noindent (6) 시계열 도면 자료(소재지: A시 B구 C동 10번지 일원) \par
\qfigure[\linewidth]{figures/03-maps.png}
8. 기타 참고자료 \par
\medskip\noindent (1) 「도시철도법」 제10조 제2항: 동법 제7조 제1항에 따른 사업계획의 승인과 같은 조 제6항에 따른 고시는 「공익사업을 위한 토지 등의 취득 및 보상에 관한 법률(이하 “토지보상법”)」 제20조 제1항 및 제22조에 따른 사업인정 및 사업인정고시로 봄 \par
\medskip\noindent (2) 대상 토지는 C동 10번지 건축허가 이전까지 상업나지 상태였으며, 건축허가를 득하는 과정에서 분필되어 현재까지 도로로 이용 중임 \par
\medskip\noindent (3) 대상 토지의 수용재결 평가액은 15,390,000원(2개 법인 평균), 이의재결 평가액은 15,770,000원(2개 법인 평균) 임 \par
\medskip\noindent (4) C동 10번지, C동 10-1번지, C동 10-2번지, C동 10-3번지는 모두 丙 소유임 \par
\medskip\noindent (5) 공시지가기준법으로 평가한 대상 토지의 평가액은 인근 거래사례의 가격수준과 적정한 균형을 이루고 있고, 그 합리성이 인정되는 것으로 봄 \par
\medskip\noindent (6) 본 사업은 도시철도 사업으로서 「토지보상법」 제70조 제5항, 동법 시행령 제38조의2 및 동법 시행령 제37조의 검토는 불필요함 \par
\medskip\noindent (7) 토지 평가단가는 천원미만은 절사할 것 \par
\medskip\noindent (8) 개별요인은 조건 간 상승식으로 산정하되, 각 조건별 비교치는 소수점 셋째자리에서 반올림하여 둘째자리까지 표시하고, 개별요인 비교치는 소수점 넷째자리에서 반올림하여 셋째자리까지 표시함 \par
\end{longquestion}
